\documentclass[11pt]{article}

% Packages
\usepackage[margin=1in]{geometry}
\usepackage{listings}
\usepackage{xcolor}
\usepackage{hyperref}

% Python code styling
\lstset{
    language=Python,
    basicstyle=\ttfamily\small,
    keywordstyle=\color{blue},
    commentstyle=\color{gray},
    stringstyle=\color{red},
    numbers=left,
    numberstyle=\tiny\color{gray},
    stepnumber=1,
    numbersep=5pt,
    backgroundcolor=\color{white},
    showspaces=false,
    showstringspaces=false,
    showtabs=false,
    frame=single,
    tabsize=4,
    captionpos=b,
    breaklines=true,
    breakatwhitespace=false,
}

\title{Sakthy re-learns ML properly}

\begin{document}

\maketitle
\tableofcontents
\newpage

\section{Motivation}
I have learned Python before, but as time has passed, I find myself needing a refresher to get back up to speed with the latest features and best practices. People say there are many ways Python can be written. But I always felt there is one true way to write Python depending on the use case. So this document serves as my internal monologue as I make these decisions throughout my workflow. It contains code snippets, explanations and occasional ramblings about why I choose to do things a certain way. Written by me for me.

\section{File Organisation \& Structure}

\subsection{What's in \texttt{\_\_init\_\_.py}?}



\end{document}
